% !PRINTED 460
\weberpage{460}
%<<CONT
die entwickelt die Form erhält
\weq{2}{\Delta = 5x_1^2x_2^2x_3^2 - x_1x_3^5 - x_2x_1^5 - x_3x_2^5.}
Auf der Curve $\Delta$ liegen also die Ecken des Coordinatendreiecks,
und folglich liegen alle siebenzähligen Pole $P_7$ auf $\Delta$ und
bilden das vollständige System der Schnittpunkte von $f$ und $\Delta$.

Eine weitere Covariante, und zwar vom 14$^{\mathrm{ten}}$ Grade,
erhalten wir aus \S.\,87, (3), nämlich die Determinante:
\weq{3}{C = \frac{1}{9}
\begin{vmatrix}
f''(x_1,x_1), & f''(x_1,x_2), & f''(x_1,x_3), & \Delta'(x_1) \\
f''(x_2,x_1), & f''(x_2,x_2), & f''(x_2,x_3), & \Delta'(x_2) \\
f''(x_3,x_1), & f''(x_3,x_2), & f''(x_3,x_3), & \Delta'(x_3) \\
\Delta'(x_1), & \Delta'(x_2), & \Delta'(x_3), & 0
\end{vmatrix}.}
Die Determinante $C$ lässt sich aus (3), wenn auch etwas weitläufig,
berechnen, und erhält den Ausdruck:
\weq{4}{C = \Sigma x_1^{14} - 34x_1x_2x_3\,\Sigma x_1^{10}x_3 - 250x_1x_2x_3\,\Sigma x_1^8x_2^3 + 375x_1^2x_2^2x_3^2\,\Sigma x_1^6x_2^2 + 18\,\Sigma x_1^7x_2^7 - 126x_1^3x_2^3x_3^3\,\Sigma x_1^3x_2^2,}
worin das Zeichen $\Sigma$ bedeutet, dass die Summe der drei Glieder
genommen werden soll, die man aus dem ersten erhält, wenn man
$x_1, x_2, x_3$ cyklisch vertauscht. Es genügt schon die Berechnung
des ersten Gliedes $x_1^{14}$, um zu erkennen, dass die Curve $C = 0$
\gesperrt{nicht} durch die Ecken des Coordinatendreiecks geht.

Eine vierte Invariante, und zwar vom 21$^{\mathrm{sten}}$ Grade,
erhalten wir nach \S.\,87, (6), wenn wir die Functionaldeterminante
der drei Formen $f$, $\Delta$, $C$ bilden:
\weq{5}{K = \frac{1}{14}
\begin{vmatrix}
f'(x_1), & \Delta'(x_1), & C'(x_1) \\
f'(x_2), & \Delta'(x_2), & C'(x_2) \\
f'(x_3), & \Delta'(x_3), & C'(x_3)
\end{vmatrix},}
die gleichfalls hieraus berechnet werden kann. Wir führen hier nur die
drei ersten Glieder an, aus denen man sieht, dass auch diese Curve
nicht durch die Ecken des Coordinatendreiecks geht
\weq{6}{K = x_1^{21} + x_2^{21} + x_3^{21} + \cdots{}^{1)}.}

$^{1)}$ Die vollständig ausgerechneten Ausdrücke finden sich in der
Abhandlung von \gesperrt{Gordan}: \glqq Ueber die typische Darstellung
der ternären biquadratischen Form $f = x_1^3x_2 + x_2^3x_3 + x_3^3x_1$\grqq{}
[Mathem. Annalen, Bd. XVII, S. 366 (1880)]. Auch in \gesperrt{Klein-Fricke},
Modulfunctionen, Bd. I, S. 734.
